% A.82 reserved by the root.
\subsection{The common multiplier and a finite quadratic-form test}

\begin{statement}{A.82}{Proved}
{A multiplier-independent quadratic obstruction recognizes the parameter thirty-four}
Let $\lambda\in2\Z\setminus\{0\}$, put $k=v_2(\lambda)$,
and write
\[
                  |1-\lambda|=\prod_{q\in Q}q^{e_q}.
\tag{A.82.1}
\]
Let $P$ be the odd primes dividing $\lambda$, together
with two when $k\ge5$.  Retain the comparison scopes of
A.79: an arbitrary labeled integral realization when
$k\notin\{2,4\}$, and full $H$-conjugacy when
$k\in\{2,4\}$.
For each $p\in P$ form the diagonal quadratic form
\[
             Q_p((Z_q)_{q\in Q})
                         =\sum_{q\in Q}e_q\log_p(q)Z_q^2.
\tag{A.82.2}
\]
A local cusp of label zero at $p$ supplies a zero of this
form with every $Z_q\in\Z_p^\times$.
In particular, when every $Q_p$, $p\in P$, has no such unit zero,
the section is exactly $s_\lambda$.

This unit-zero question has the following exact finite test.
Put $d_p=\min_{q\in Q}w_{p,q}$, with the explicit integers
$w_{p,q}$ of (A.79.3).
For odd $p$ put
\[
 a_{p,q}=\frac{e_q(q^{p-1}-1)}{p^{d_p}}\pmod p,
 \qquad
 T_p=\left\{\sum_{q\in Q}a_{p,q}z_q^2:
                                      z_q\in\F_p^\times\right\}.
\tag{A.82.3}
\]
The indicated integer quotient is defined because
$w_{p,q}\ge d_p$.
Then (A.82.2) has a unit zero if and only if $0\in T_p$.
When $p=2$, it has a unit zero if and only if
\[
                                  k\ge d_2+3.
\tag{A.82.4}
\]
For two prime factors $Q=\{q_1,q_2\}$, the equivalent
description is equality of the two weights and
\[
             -\frac{e_{q_1}\log_p(q_1)}
                      {e_{q_2}\log_p(q_2)}
                            \in(\Z_p^\times)^2.
\tag{A.82.5}
\]
At odd $p$ this is a nonzero square-residue test modulo $p$;
at two its unit residue must equal one modulo eight.

The parameter $\lambda=34$ is recognized with an arbitrary
labeled integral comparison.  It is the first positive
integer not included in the previously listed parameter
ranges A.41, A.45, A.63, A.64 and A.79.
The last observation refers to these explicit ranges,
not to a classification of every earlier recognition method.
\end{statement}
\begin{proof}
The support argument of A.79 also applies to negative even
integers.  Odd divisors of $\lambda$ have contact zero,
odd divisors of $1-\lambda$ have contact one, and there
are no infinity cusps.  At two the same valuation parity,
exact ordinary-value and, in the two exceptional cases,
mod-eight arguments give $v_2(B)=0$.
Only the global sign changes.  Thus
\[
          B=\operatorname{sgn}(1-\lambda)
                                  \prod_{q\in Q}q^{\beta_q}.
\tag{A.82.6}
\]

We retain the multiplier rather than replace the given
comparison by a symplectic one.  Fix, once for all places,
the global integral intertwiners
$g_\ell:T_\ell(E_\lambda)\xrightarrow{\sim}
T_\ell(E_U)|_s$ of A.16 and put
$\Delta_\ell=\det(g_\ell)\in\Z_\ell^\times$, using
the fixed Weil-paired geometric bases.
At a local prime $q\in Q$, the representative change
from the localized section to its point or cusp section
is geometric; its Legendre matrix lies in $H$ and has
determinant one.  Passing to a common finite splitting
extension and to paired Tate bases also preserves the
Weil pairings and hence the multiplier.
Therefore every such local triangular comparison has
diagonal entries $a_{q,\ell},d_{q,\ell}$ satisfying
\[
 a_{q,\ell}d_{q,\ell}=\Delta_\ell,
 \qquad
 \frac{\beta_{q,\ell}}{e_q}
       =\frac{a_{q,\ell}}{d_{q,\ell}}
       =\frac{a_{q,\ell}^{2}}{\Delta_\ell},
 \qquad a_{q,\ell},d_{q,\ell}\in\Z_\ell^\times.
\tag{A.82.7}
\]
For a cusp this is the full Tate-extension valuation map
for $B^2/256$.  For an ordinary local point, (A.78.1)
gives $\beta_q=e_q$ and the same calculation gives
$a_{q,\ell}=d_{q,\ell}$; thus (A.82.7) holds there too.
The quantity $\Delta_\ell$ is the same for every $q$;
the paired basis changes have not introduced separate
local multipliers.

At a label-zero cusp $p$, regularity gives $B_p=1$.
The coefficient-$p$ logarithm of (A.82.6) is
$\sum_q\beta_{q,p}\log_p(q)=0$.
Substitute (A.82.7) and multiply by the single nonzero
unit $\Delta_p$.  The result is exactly
\[
                   \sum_q e_q\log_p(q)a_{q,p}^2=0,
                       \qquad a_{q,p}\in\Z_p^\times.
\tag{A.82.8}
\]
This proves the necessary unit-zero condition without
asserting that the separate $\beta_{q,p}/e_q$ are squares.
When the test excludes every label-zero cusp, only
label one remains and A.37 gives the point conclusion;
the labeled primitive-isogeny argument fixes its exact
parameter $\lambda$.

We prove both directions of the finite test.
Write $c_q=e_q\log_p(q)$ and divide the form by
$p^{d_p}$.  Its coefficients are integral and at least
one, say $c_{q_0}/p^{d_p}$, is a unit.
For odd $p$, the logarithm series and (A.79.8) give
\[
 \frac{c_q}{p^{d_p}}
       \equiv\frac{a_{p,q}}{p-1}\pmod p.
\tag{A.82.9}
\]
Terms beyond the first in $\log(q^{p-1})$ have valuation
at least one more than its first term; multiplication
by $e_q$ preserves the required bound after normalization.
Thus reduction of the form is, up to one common nonzero
factor, the sum in (A.82.3).
A unit zero reduces to a tuple in $(\F_p^\times)^Q$.
Conversely choose such a zero modulo $p$ and arbitrary
unit lifts.  With every variable but $Z_{q_0}$ fixed,
the derivative is $2(c_{q_0}/p^{d_p})Z_{q_0}$, a unit.
Hensel's lemma gives an actual zero with the same unit
residues.  This proves the stated equivalence at odd $p$.

At two, every unit square equals one modulo eight.
A unit zero therefore requires
\[
                        \sum_q c_q/2^{d_2}=0\pmod8.
\tag{A.82.10}
\]
Conversely, if (A.82.10) holds, start all variables at one
and vary only the variable with unit coefficient.
The polynomial value has valuation at least three and
its derivative has valuation one.  The strict Hensel
inequality $3>2$ gives a root congruent to one modulo four,
so every variable remains a unit.
Now $\sum_q c_q=\log_2|1-\lambda|=\log_2(1-\lambda)$,
since the sign has zero logarithm.  As $p=2$ belongs to
$P$ only when $k\ge5$, this logarithm has exact valuation
$k$.  Condition (A.82.10) is exactly $k\ge d_2+3$.
This proves (A.82.4), including its sufficiency for the
unit quadratic form.  It asserts no lift from this
quadratic-form solution to a full section.
For two variables, dividing by the nonzero second
coefficient and by the square of the first unit gives
(A.82.5).  The exact odd-prime and dyadic square-unit
criteria follow respectively from simple-root Hensel
lifting and $(\Z_2^\times)^2=1+8\Z_2$.

For $\lambda=34$, $k=1$, $P=\{17\}$ and
$Q=\{3,11\}$, with both exponents one.
The exact residues are
\[
 \frac{3^{16}-1}{17}\equiv10,\qquad
 \frac{11^{16}-1}{17}\equiv9\pmod{17}.
\tag{A.82.11}
\]
Both weights equal one.  A unit zero would force
\[
                       (Z_{11}/Z_3)^2
                               =-10/9=14\pmod{17}.
\tag{A.82.12}
\]
The nonzero squares modulo seventeen are
$1,2,4,8,9,13,15,16$, so fourteen is excluded.
The general criterion recognizes $34$ with no multiplier
restriction.

We also record the exact coefficient-two and
coefficient-three constraints, to expose rather than
discard the possible dyadic alternative in a direct
exponent calculation.
The dyadic local point is potentially good.
The valuation parities from A.16 allow
$(v_2(A),v_2(B))=(1,0)$ or $(3,0)$.
The second alternative would make $1-x_2$ a square,
since it is one modulo eight.  But its prescribed
squareclass is $-33\equiv7\pmod8$, which is not a
square.  Hence the exact global support has the form
\[
                  A=2\cdot17^\alpha,
               \qquad B=-3^\beta11^\gamma.
\tag{A.82.13}
\]
If the section were cuspidal at three, its label would
be one and its regular coordinate would impose
$\log_3(2)+\alpha_3\log_3(17)=0$.
But
\[
              v_3(\log_3 2)=1,
                       \qquad v_3(\log_3 17)=2,
\tag{A.82.14}
\]
whereas $\alpha_3$ is a unit by the integral Tate
valuation comparison at seventeen.  This is impossible.
Thus three is ordinary, its exact value is $34$ by
A.78, and A.34 makes the global multiplier a square.
In this resulting full-$H$ comparison a putative
label-one cusp at eleven would require
\[
              \alpha_{11}
                   =-\frac{\log_{11}2}{\log_{11}17}
                         \equiv-5/3=2\pmod{11}.
\tag{A.82.15}
\]
Here the normalized Fermat quotients are five and
three.  The number two is a nonsquare modulo eleven,
whereas (A.82.7) now makes $\alpha_{11}$ a square unit.
This gives a second complete recognition route with
the exact dyadic and coefficient-three branches retained.

The range comparison is literal: A.45 treats $2$ through
$31$, $32$ satisfies A.41 and A.63, and $33$ satisfies
A.63.  The parameter $34$ has a composite odd neighbour,
neither consecutive integer is dyadic, its dyadic
valuation is one, and its two A.79 weights at seventeen
are equal.  Therefore none of the listed tests supplies
it before (A.82.12).
\end{proof}

\noindent The exact arithmetic and finite quadratic-form
certificate is \path{code/legendre_square_form_certificate.py},
with retained output
\path{code/legendre_square_form_certificate_output.txt}.
It checks the residue forms, the dyadic unit-square test,
and every displayed numerical obstruction.  The proof
retains the common global multiplier throughout; it
does not assert a solution of the general recognition
problem.
